Полученные выводы следует интерпретировать с учётом ограничений
бенчмарка. Для каждой комбинации задачи, метода и бюджета использовано
только 10 целевых наблюдений. Ограничение связано с вычислительной
стоимостью экспериментов; увеличение числа posterior-сэмплов для каждого
наблюдения не заменяет увеличение числа независимых целевых наблюдений.
Поэтому агрегированные показатели калибровки, включая TARP и SBC, имеют
ограниченную статистическую мощность и носят ориентировочный характер.

Для SIR и StochVol нет аналитического эталонного posterior, поэтому
качество оценивалось по калибровке и posterior-predictive проверкам.
Хорошая калибровка сама по себе не означает, что оценённая плотность
совпадает с истинным posterior. Кроме того, в StochVol в бенчмарке
используется пятикомпонентная сводка доходностей вместо полного временного
ряда; выводы относятся к posterior параметров при условии этой сводки.

Отдельное ограничение связано с поддержкой распределений. Для части
сэмплов SFMPE и SMC-ABC в StochVol нарушаются условия $-1<\phi<1$ и
$\sigma>0$; следовательно, соответствующие оценки калибровки нужно
рассматривать вместе с долей сэмплов, попавших в допустимую область.
Для SFMPE эта доля составляет приблизительно $0.33$, $0.36$ и $0.77$
при бюджетах 100, 1000 и 10000 соответственно; для SMC-ABC --- $0.61$,
$0.76$ и $0.78$. Поэтому близкий к нулю ATC SFMPE на малых бюджетах
нельзя интерпретировать отдельно от доли сэмплов вне допустимой области.

Наконец, измеренное время выполнения зависит от аппаратного устройства
и фоновой нагрузки, а запуски бенчмарка выполнялись с автоматическим
выбором устройства. Поэтому приведённые значения времени дают лишь
ориентировочное сравнение. Бюджет симуляций учитывает симуляции,
использованные для инференса, но не posterior-сэмплирование и
дополнительные симуляции для диагностик.